В ходе выполнения курсовой работы была рассмотрена задача интерактивной визуализации двумерного течения жидкости с использованием современных веб технологий. Были изучены физико-математические основы моделирования гидро- и газодинамических процессов, а также некоторые алгоритмы визуализации модели жидкости. На основе исследования был выбран и реализован алгоритм Stable Fluids, который обеспечивает баланс между качеством визуализации, вычислительной эффективностью и сложностью реализации.


В работе был использован подход к реализации симуляции жидкости на графическом процессоре с использованием вычислительных шейдеров. Это позволило эффективно обрабатывать большие массивы данных, описывающие поле скоростей и скалярные величины, и тем самым достичь большей производительности. Было уделено внимание правильной настройке графического и вычислительного конвейера, организации данных в виде буферов, и разделению этапов симуляции на отдельные, вычислительные шаги, которые следовали алгоритму Stable Fluids.


Результатом работы стало интерактивное веб-приложение, демонстрирующее визуализацию течения жидкости в реальном времени и позволяющее пользователю взаимодействовать с симуляцией. Это подтверждает, что современные браузерные технологии позволяют эффективно взаимодействовать с GPU и решать задачи визуализации, которые ранее были бы реализовывали только в качестве нативных приложений.

% В дальнейшем развитие данной работы может быть направлено на расширение физической модели, повышение численной устойчивости и качества визуализации, а также на исследование трёхмерных течений и более сложных граничных условий. Кроме того, было бы интересно сравнить производительности WebGPU с настольными GPU-решениями (например чистый Vulkan) на более сложных задачах симуляции жидкости.
